\documentclass[12pt]{article}

% Packages
\usepackage[utf8]{inputenc}
\usepackage[T1]{fontenc}
\usepackage{amsmath}
\usepackage{amssymb}
\usepackage{amsthm}
\input{common-defs}
\usepackage{graphicx}
\usepackage{natbib}
\usepackage{enumitem}
\usepackage{booktabs}        % For \toprule, \midrule, \bottomrule in tables
\usepackage{threeparttable}  % For table notes
\input{typography-preamble}

% Note: all project macros (\ehat, \mhat, \what, \that, \thetahat, \pihat,
% \Vhat, \Ltwo, \Linf, \argmin, \argmax, \Var, \abs, \card, theorem/lemma/
% definition/assumption environments, \E, \Prob, \cX, \cT, \cS, etc.) come
% from common-defs.tex above. Do NOT redefine them here -- that is what
% broke compilation (theory.tex's preamble was pasted in verbatim; it is a
% standalone "portable preamble" with no project macros, meant to duplicate
% common-defs.tex, not sit alongside it).

% Title information
\title{Interpretable causal inference for high-stakes settings}
\author{Denis Agniel}
\date{\today}
\begin{document}

\maketitle
\begin{abstract}
Modern causal inference relies on flexible, black-box nuisance estimation together with sample splitting to obtain valid, efficient inference for treatment effects. In high-stakes settings -- health care payment, treatment triage, sentencing -- that estimation step is itself the object stakeholders need to audit, and a black-box nuisance model, however well it performs, does not offer an auditable trail. We propose \emph{doubletree}, an estimator of the average treatment effect on the treated (ATT) that replaces the propensity and outcome nuisances with two globally optimal, size-budgeted decision trees, each small enough to display and each fit on every observation, with no sample splitting or cross-fitting. When the two nuisances are exactly representable by trees within the analyst's budget -- a condition we make precise and verifiable (Condition~(P)) -- doubletree attains the classical semiparametric efficiency bound for the ATT at the parametric $\sqrt n$ rate, and its influence-function variance estimator is the ordinary plug-in, with no jackknife or debiasing correction required. When the budget-constrained class does not contain the truth, we do not assume the approximation error away: an anchor-interval construction absorbs the resulting bias into an honest, approximation-aware confidence interval, so that a defensible inferential statement remains available even outside the class. We are explicit about what this buys and does not: the oracle property that drives our results is pointwise, not uniform, in the underlying data-generating process, and we make no claim of matching the flexibility of an unrestricted machine-learning ensemble. The comparison we defend is against a generalized linear working model: doubletree admits richer, non-additive structure than a GLM while remaining exactly as auditable, and it delivers sharper, full-sample inference than generic double machine learning whenever its structural class holds.
\end{abstract}

\section{Introduction}

Modern causal inference has focused productively on estimators that remain valid under essentially arbitrary nuisance estimation. Augmented inverse-probability weighting, double machine learning, and targeted maximum likelihood estimation are all known to deliver $\sqrt n$-consistent, asymptotically normal estimates of a treatment effect when the propensity score and outcome regression are estimated by almost any sufficiently accurate method, provided the resulting nuisance estimates satisfy standard product-rate conditions. This robustness to the nuisance-estimation method is what licenses the now-standard practice of pairing these estimators with black-box learners such as random forests or gradient boosting, chosen because they can plausibly be assumed to approximate the true nuisance functions under only mild smoothness conditions. In this division of labor the nuisance model is a means to an end -- a \emph{nuisance} in the technical as well as the colloquial sense -- and is neither expected nor required to be interpreted on its own terms.

That division of labor is unworkable in settings where the stakes of the decision demand that the estimation itself be auditable. If a causal effect estimate is used to set a health care payment adjustment or to inform a sentencing or triage decision, stakeholders may reasonably require that the calculation be traceable step by step, not merely validated by its asymptotic properties. Existing approaches to flexible causal effect estimation do not offer this: they can deliver flexible nuisance estimation with asymptotic guarantees, or they can deliver an interpretable nuisance model, but not both together with full-sample inference. Linear working models are one interpretable option, and decision trees are another, encoding a different and often more natural kind of structure -- non-additive, region-varying dependence -- than a linear model can. Section~\ref{sec:lasso-comparison} compares the assumptions under which penalized linear and tree-based nuisance estimators support valid and efficient inference. Tree-based nuisance estimators have been incorporated into causal effect estimation before, but existing constructions require sample splitting to obtain valid inference -- which makes the displayed tree an incomplete audit trail, since part of the estimate comes from a fold the reader cannot inspect against the tree -- or they sacrifice either interpretability or an asymptotic guarantee to avoid it.

We propose \emph{doubletree}, an approach to interpretable treatment effect estimation for the average treatment effect on the treated that offers flexible, data-adaptive nuisance estimation, full interpretability of every step, and asymptotic guarantees together, with no sample splitting anywhere. The name reflects the estimator's structure: one globally optimal decision tree for the propensity score and one for the outcome function, each built to a leaf budget the analyst sets in advance and each small enough, in the intended regime, to fit on an index card. The budget is a single, user-specified interpretability knob: a smaller budget gives a more displayable tree, at a cost we quantify precisely rather than leave implicit. Under conditions that make the resulting tree partitions exactly recover the population structure of the two nuisances -- Condition~(P), stated formally in \S\ref{sec:partition-recovery} -- doubletree's estimate converges to the true treatment effect at the parametric rate, with no correction needed for having chosen the tree structure adaptively from the same data used for inference. When the true nuisances are not exactly representable within the leaf budget, doubletree does not silently degrade: \S\ref{sec:honest-manuscript} constructs an anchor-interval confidence set that remains valid by absorbing the unavoidable approximation bias into its width rather than assuming it away. None of this requires reserving any observations for a separate selection or evaluation step; every nuisance is fit, and every inferential guarantee holds, using the entire sample.

The remainder of the paper proceeds as follows. Section~\ref{sec:setup} sets up notation, the tree partition class, and the doubletree algorithm itself. Section~\ref{sec:instantiation1} (fixed grid, exact sparsity) and \S\ref{sec:instantiation2} (continuum sparsity with off-grid threshold refinement) each instantiate Condition~(P) under progressively weaker restrictions on where the analyst's candidate splits may fall. Section~\ref{sec:honest-manuscript} addresses what happens, and what can still be claimed, when Condition~(P) fails. Section~\ref{sec:lasso-comparison} situates doubletree against penalized linear methods, and closes by stating plainly what doubletree does not claim: the recovery guarantee underlying our results holds pointwise in the data-generating process, not uniformly as the signal defining the true tree structure shrinks, and we do not claim to match the flexibility of an unrestricted machine-learning ensemble -- only that of a generalized linear working model, with interpretability and full-sample inference held fixed.


\section{Setup, identification, and algorithm}
\label{sec:setup}


\subsection{Data, estimand, and notation}

We observe $n$ observations $O_1, \ldots, O_n$ each an \iid draw from the random variable $\bO = (\bX, A, Y) \sim \Prob$. The observed data $\bO$ consists of a covariate vector $\bX \in \cX \subset \R^p$, a binary treatment $A \in \{0,1\}$, and an outcome $Y \in \cY \subset \R$. Let $Y(a)$ denote the potential outcome under treatment $A = a$. Define the two primary nuisance functions to be the propensity score $e_0(\bx)=\Prob(A=1\mid \bX=\bx)$ and the control ($A=0$) outcome function $\mu_0(\bx)=\E[Y\mid A=0,\bX=\bx]$. Define also $\pi=\Prob(A=1)$, and let $m_1(\bx)=\E[Y\mid A=1,\bX=\bx]$ and $\sigma_a^2(\bx)=\Var(Y\mid A=a,\bX=\bx)$. 

Our target parameter is the average treatment effect on the treated (ATT),
\begin{equation}
\label{eq:att}
\theta_0 = \E[Y(1) - Y(0) \mid A=1].
\end{equation}

It is useful to define the following weight when working with the ATT: $w_0(\bx)=e_0(\bx)/\{1-e_0(\bx)\}$, and for a candidate propensity $e(\bx)$ we write
$w(\bx)=e(\bx)/\{1-e(\bx)\}$. We collect the two nuisances into $\eta_0=(e_0,\mu_0)$, and define candidate nuisances as $\eta=(e,\mu)$. It is convenient to index the two nuisances by $j\in\{e,\mu\}$ and to
write $\gamma_{0,e}=e_0$, $\gamma_{0,\mu}=\mu_0$.

We write $\Prob_n$ for the empirical measure, $\Prob_n f=n^{-1}\sum_{i\le n}f(\bO_i)$, and
$\pihat=\Prob_n A$. Two norms on functions of $\bx$ are used throughout:
\[
  \twonorm{f}^2=\E\bigl[f(\bX)^2\bigr],
  \qquad
  \|f\|_{2,w}^2=\E\bigl[\{1-e_0(\bX)\}f(\bX)^2\bigr],
\]
the second being the control-weighted norm. Convergence in distribution is
denoted $\rightsquigarrow$, and $z_{1-\alpha/2}$ is the standard normal
quantile.


\subsection{Covariate space, tree partitions, and optimal trees}\label{sec:trees}

The covariate space $\cX$ need not be finite. Instead, before seeing the data,
the analyst fixes a finite collection of permissible cutpoints or category
groupings on each coordinate, such as age brackets, income bands, or clinical
cutoffs. These cutpoints are never updated using the observed data and generate
a finite collection of grid atoms $\cA$, fixed in $n$, with
\[
  M:=\card{\cA}<\infty,
  \qquad
  \Prob(X\in a)>0\quad\text{for every }a\in\cA.
\]
A \emph{tree partition} $\tau$ of $\cX$ is induced by recursive axis-aligned
binary splits at these pre-specified cutpoints only. Its blocks are
\emph{leaves}, each leaf is a union of grid atoms, $\card\tau$ is its leaf
count, and $\ell_\tau(x)$ is the leaf containing $x$. This restriction on the
permissible splits is what makes the candidate partition class fixed and
finite; it is not merely a convention. Further, for leaf $\ell$ define the leaf sample size as $|\ell| = \sum_{i=1}^n \cI\{X_i \in \ell\}$, and the control sample size as $|\ell|_0 = \cI\{X_i \in \ell, A_i = 0\}$. For a strictly positive weight $\nu$ on $\cX$, let $\Pi^\nu_\tau$ denote the $\nu$-weighted projection onto functions constant on the leaves of $\tau$; we write $\Pi_\tau$ for $\Pi^{p}_\tau$, and use the control-weighted version, $\nu_x=p_x\{1-e_0(x)\}$, for the outcome nuisance. Because every leaf is a nonempty union of positive-mass atoms, every candidate leaf automatically has positive probability mass; no separate raw-cell probability floor is required.

We operate under the assumption that a tree is sufficiently interpretable if it obeys a user-specified \emph{leaf budget}, or a maximum number of a leaves the tree partition may have. This budget allows the user to determine how complex the tree may become. The saturated grid budget is $\bar L = M$, one leaf per grid atom. In general, stakeholders may prefer a smaller tree, in which case we set the leaf budget at $\bar L \leq M$. Then, define the space of all possible partitions under that budget as
\[
  \cT_{\bar L}=\bigl\{\tau:\tau\text{ a tree partition of }\cX\text{ built
  from the pre-specified cutpoints},\
  \card\tau\le\bar L\bigr\}.
\]
Since the cutpoints and grid atoms are fixed and finite, $\cT_{\bar L}$ is also fixed and finite.

Finally, we define for any function $\gamma$ a \emph{sufficient class}, or a class of trees that are sufficient to describe the function comprehensively:
 \begin{definition}[Sufficient class]\label{def:sufficient}
For $j\in\{e,\mu\}$,
\[
  \cS_j=\bigl\{\tau\in\cT_{\Lbar}:
  \Pi^\nu_\tau\gamma_{0,j}=\gamma_{0,j}\ \ \Prob\text{-a.s.}\bigr\},
\]
equivalently the set of $\tau\in\cT_{\Lbar}$ on whose leaves $\gamma_{0,j}$ is
$\Prob$-a.s.\ constant. The characterisation is independent of the strictly
positive weight $\nu$.
\end{definition}
 
 \subsection{Identification and influence function}
 
In this section, we briefly review how the causal effect of interest is identified. 

\begin{assumption}[Causal identification]\label{ass:causal}
Consistency, $Y=AY(1)+(1-A)Y(0)$; ignorability,
$\{Y(0),Y(1)\}\perp A\mid X$; $\pi>0$; and there exists $c\in(0,1)$ with
\[
  1-e_0(x)\ \ge\ c \qquad\text{for every }x\in\cX .
\]
\end{assumption}

Versions of this set of assumptions are standard in the causal inference literature, though we require a \emph{uniform} lower bound $1-e_0(x)\ge c$ for every $x\in\cX$, rather than the weaker pointwise condition $e_0(x)<1$ that suffices in some other treatments. On a finite covariate space the two are essentially interchangeable, since a maximum over finitely many finite values is automatically finite. On a general (e.g.\ continuous) $\cX$ they are not: the semiparametric efficiency bound's variance decomposition (the companion theory document's Eq.~(vdecomp)) contains a term $\E[e_0(X)^2\sigma_0^2(X)/\{1-e_0(X)\}]$, which can diverge under only pointwise positivity if $P_X$ places mass where $1-e_0(x)\to0$ without a compensating decay in the control-outcome variance $\sigma_0^2(x)$. The uniform bound $c$ is what keeps this term finite, not merely a convenience. Under Assumption \ref{ass:causal}, we may identify the ATT by the following observed data quantity:
\begin{align}
	\theta_0 = \E\left\{Y - \E\left(Y | \bX, A = 0\right) | A = 1\right\}.
\end{align}
The efficient influence function for this functional in the nonparametric model is
\citep{hahn1998}
\begin{equation}\label{eq:score}
  \psi(O;\theta,\eta)
  =\frac1\pi\Bigl[A\bigl\{Y-\mu(X)-\theta\bigr\}
   -(1-A)\,w(X)\bigl\{Y-\mu(X)\bigr\}\Bigr],
  \qquad w=\frac e{1-e},
\end{equation}
and we write $\psi_0=\psi(\cdot;\theta_0,\eta_0)$. We will use this quantity to build a new estimator for the ATT based on interpretable trees in the next section.

\subsection{The doubletree algorithm}
To build an intepretable causal estimate, we use globally optimal intepretable trees to estimate both required nuisance functions $e_0, \mu_0$, and -- since this approach relies on two trees -- we call our estimation approach the \emph{doubletree} algorithm.

Doubletree select the tree from $\cT_{\bar L}$ that minimizes the penalised empirical risk over the space of feasible trees. We take the space of feasible trees to be a subset of $\cT_{\bar L}$ for regularity reasons. Specifically, define the feasible sets as
\begin{equation}\label{eq:feasible}
  \cT^{(\mu)}_n=\bigl\{\tau\in\cT_{\bar L}: \min_{\ell \in \tau} |\ell|_0 \geq m_n\bigr\}, \qquad 
    \cT^{(e)}_n=\bigl\{\tau\in\cT_{\bar L}: \min_{\ell \in \tau} |\ell| \geq m_n\bigr\},
\end{equation}
and for a penalty $\lambda_n>0$ set
\begin{equation}\label{eq:select}
  \that_j\in\argmin_{\tau\in\cT^{(j)}_n}
  \bigl\{R^{(j)}_n(\tau)+\lambda_n\card\tau\bigr\},
  \qquad j\in\{e,\mu\},
\end{equation}
where $R^{(j)}_n(\tau)$ is the empirical risk of tree $\tau$ for nuisance $j$,
\begin{equation}\label{eq:risk-emp}
R^{(e)}_n(\tau)=\Prob_n\bigl[\{A-\hat e_\tau(\bX)\}^2\bigr],
\qquad
R^{(\mu)}_n(\tau)=\frac{1}{n_0}\sum_{i:A_i=0}\bigl\{Y_i-\hat\mu_\tau(\bX_i)\bigr\}^2,
\end{equation}
squared error for both the propensity tree and the outcome tree, with $n_0=\sum_{i\le n}(1-A_i)$ and leaf-wise empirical plug-ins
\begin{equation}\label{eq:refit}
\hat e_\tau(\bx)=\frac{1}{|\ell_\tau(\bx)|}\sum_{i:\bX_i\in\ell_\tau(\bx)}A_i,
\qquad
\hat\mu_\tau(\bx)=\frac{1}{|\ell_\tau(\bx)|_0}\sum_{i:\bX_i\in\ell_\tau(\bx),\,A_i=0}Y_i ,
\end{equation}
i.e.\ the empirical proportion treated and the empirical control-outcome mean within the leaf containing $\bx$. Given leaf-refit nuisances $(\hat e_{\hat\tau_e},\hat\mu_{\hat\tau_\mu})$ for the selected partitions $\hat\tau_e,\hat\tau_\mu$, the doubletree estimate $\thetahat=\thetahat(\hat\tau_e,\hat\tau_\mu)$ is the solution of the empirical estimating equation $\Prob_n[\psi(\bO;\theta,\hat e_{\hat\tau_e},\hat\mu_{\hat\tau_\mu})]=0$ for $\psi$ as in \eqref{eq:score}, i.e.\
\[
  \thetahat
  =
  \frac{1}{n\hat\pi}\sum_{i=1}^n
  \Bigl[A_i\{Y_i-\hat\mu_{\hat\tau_\mu}(\bX_i)\}
    -\{1-A_i\}\hat w_{\hat\tau_e}(\bX_i)\{Y_i-\hat\mu_{\hat\tau_\mu}(\bX_i)\}\Bigr],
\]
with $\hat\pi=\Prob_n A$ and $\hat w_{\hat\tau_e}=\hat e_{\hat\tau_e}/(1-\hat e_{\hat\tau_e})$; the root is unique whenever $\hat\pi>0$, since $\psi$ is affine and strictly monotone in $\theta$. The quantity $R^{(j)}(\tau)$ is the population counterpart of $R^{(j)}_n(\tau)$,
\begin{equation}\label{eq:risk-pop}
R^{(e)}(\tau)=\E\bigl[\{A-\Pi_\tau e_0(\bX)\}^2\bigr],
\qquad
R^{(\mu)}(\tau)=\E\bigl[\{Y-\Pi^\nu_\tau\mu_0(\bX)\}^2\mid A=0\bigr],
\end{equation}
obtained from \eqref{eq:risk-emp} by replacing the empirical measure with the population measure and the empirical leaf plug-ins $\hat e_\tau,\hat\mu_\tau$ with the projections $\Pi_\tau e_0$ and $\Pi^\nu_\tau\mu_0$ introduced above. These projections are exactly the squared-error minimizers.

\begin{remark}
Note that we are using a flexible tree learner to estimate both nuisance functions, but we do not require sample-splitting/cross-fitting as most approaches that use data-adaptive nuisance learners do \citep{chernozhukov2018}. This means that doubletree will make more efficient use of potentially limited data. It also means that doubletree is fully explained by exactly two trees. Under $K$-fold cross-fitting, users would need to interpret 2$K$ trees, as opposed to doubletree's two.
\end{remark}

\begin{remark}[Why two trees: interpretability and identification are separate questions]
\label{rem:two-trees-why}
The account above motivates two trees on interpretability grounds -- doubletree is exactly displayable as two objects -- but interpretability alone does not explain why \emph{efficiency} or \emph{validity} requires both nuisances to be estimated by a tree. Two points are worth separating precisely. First, the propensity \emph{tree} specifically is an interpretability choice, not by itself a validity requirement: Proposition~\ref{prop:bilinear}'s bilinear structure only requires the \emph{product} of the two realized nuisance errors, $D_wD_\mu$ in the notation of Lemma~\ref{lem:biasbound}, to vanish fast enough, and that product-rate condition is indifferent to whether $\hat w$ is itself piecewise constant on a tree. This is necessary but not sufficient on its own: because doubletree performs no sample splitting or cross-fitting, asymptotic linearity of $\thetahat$ also requires the own-observation empirical-process term $(\Prob_n-\Prob)\{\psi(\cdot;\theta_0,\hat\eta)-\psi_0\}$ to be $o_p(n^{-1/2})$, and it is exactly to control \emph{that} term without splitting that Condition~(P) and the fixed, finite candidate class $\cT_{\Lbar}$ are needed (\S\ref{sec:partition-recovery}). An arbitrary non-tree $\hat w$ fit on the same $n$ observations and satisfying only the product-rate condition is not in general covered by this argument. Second, the propensity \emph{nuisance} itself is genuinely informative for the ATT, unlike for the ATE: \citet{hahn1998} shows that knowledge of $e_0$ does not lower the semiparametric efficiency bound for the average treatment effect, but it does lower it for the average treatment effect on the treated -- so dropping the propensity is not a free simplification for the estimand considered here. Corollary~\ref{cor:single-tree} and its consequences below make both points precise in the concrete special case of a single outcome tree fit with no propensity tree at all.
\end{remark}

Additionally, we assume that the fitted propensities agree with the positivity assumption in \ref{ass:causal}:
\begin{assumption}[Clipping]\label{ass:construct}
Fitted propensities are clipped so that $1-\hat e(\bx)\ge c$ for every $\bx$, with $c$ the constant of Assumption~\ref{ass:causal}.
\end{assumption}

\section{Doubletree bias and inference}

\subsection{Partition recovery and the central limit theorem}
\label{sec:partition-recovery}

Doubletree's efficiency rests entirely on whether the two fitted tree
partitions recover the population structure of the two nuisances -- not on
any property of the grid used to search for them. We isolate exactly this
requirement, state the central limit theorem once under it, and then verify
it in two settings that differ only in \emph{how} recovery is established:
first when the truth is exactly representable on the analyst's own
pre-specified grid (\S\ref{sec:instantiation1}), and second when it is
exactly representable by some tree of the same size, at threshold values not
confined to that grid (\S\ref{sec:instantiation2}).

\begin{definition}[Partition recovery -- Condition~(P)]
\label{cond:P}
For $j\in\{e,\mu\}$, let $\tauhat_j$ denote the fitted partition (topology
and threshold vector), let $\nu_e=\Prob_X$, and let
$d\nu_\mu=(1-e_0)\,d\Prob_X$. Condition~(P) holds for nuisance $j$ if:
\begin{enumerate}[label=(P\arabic*),leftmargin=*]
\item[(P0)] There is a fixed reference partition $\tau_{0,j}$ with at most
$\Lbar$ leaves on which $\gamma_{0,j}$ is $\Prob$-a.s.\ constant, with every
leaf of positive $\nu_j$-mass;
\item[(P1)] with probability tending to one, $\tauhat_j$ and $\tau_{0,j}$
have the same tree topology and the same split coordinate at each matched
internal node, so their leaves admit a canonical bijection preserving
matched split boundaries and left--right incidence. Threshold values may
differ under this correspondence, and equality of the induced leaf sets is
\emph{not} asserted -- only the existence of the bijection. Literal rooted-
or ordered-tree equality is not required;
\item[(P2)] under the leaf correspondence of (P1), let $S_{j,n}$ be the set
of observations whose leaf assignment under $\tauhat_j$ differs from their
assignment under $\tau_{0,j}$ (nonempty in general, since the matched
leaves need not coincide as sets when thresholds differ). On the event in
(P1),
$\card{S_{j,n}}=O_p(1)$.
\end{enumerate}
\end{definition}

Because the leaf values \eqref{eq:refit} are, by construction, leaf-wise
empirical means, pinning down the partition pins down the fitted nuisance:
Condition~(P) is a condition on partitions, not directly on the fitted
functions.

\begin{remark}[Why (P2) is a bounded count, not merely a vanishing rate]
\label{rem:P2-mechanism}
(P2) is stated in the currency the central limit theorem's proof actually
consumes -- a bounded number of reassigned observations -- rather than as a
threshold-rate condition, because the two are \emph{not} interchangeable
under only a finite (not bounded) second moment on $Y$. If the covariate
law has a density bounded away from zero and infinity near each true
boundary, $\card{S_{j,n}}=O_p(1)$ is equivalent to the threshold rate
$\|\hat c_j-c_{0,j}\|=O_p(n^{-1})$; a weaker rate such as
$\|\hat c_j-c_{0,j}\|=o_p(n^{-1/2})$ does \emph{not} suffice in general, since
it permits $\card{S_{j,n}}$ to grow, and the maximal outcome value among a
growing, adversarially-selected set of reassigned observations is controlled
only by the same soft $o_p(\sqrt n)$ envelope that bounds the maximum over
the \emph{entire} sample -- too weak once multiplied by a growing count. A
bounded outcome, or cross-fitting the partition-selection step away from the
leaf refits, would each independently relax (P2) to the weaker rate; neither
is assumed here, since doubletree uses neither.
\end{remark}


\begin{theorem}[Central limit theorem under partition recovery]
\label{thm:main}
Suppose the observations are \iid; Assumptions~\ref{ass:causal} and
\ref{ass:construct} hold; the conditional second moments of $Y$ are
uniformly bounded; $\Lbar$ is fixed;
\[
  V:=\E[\psi_0(\bO)^2]>0;
\]
and Condition~(P) holds for both $j=e,\mu$. Then, with both tree structures
and leaf values fitted using all $n$ observations,
\[
  \sqrt n(\thetahat-\theta_0)\rightsquigarrow N(0,V).
\]
\end{theorem}

\begin{corollary}[Variance estimation and confidence intervals]
\label{cor:variance}
Under the conditions of Theorem~\ref{thm:main}, let
\[
  \Vhat
  =
  \frac{1}{n\pihat^2}
  \sum_{i=1}^n
  \left[
    A_i\{Y_i-\hat\mu(\bX_i)-\thetahat\}
    -
    (1-A_i)\hat w(\bX_i)\{Y_i-\hat\mu(\bX_i)\}
  \right]^2,
  \qquad
  \hat w=\frac{\hat e}{1-\hat e}.
\]
Then $\Vhat\to_p V$, and
\[
  \thetahat
  \pm
  z_{1-\alpha/2}\sqrt{\Vhat/n}
\]
has pointwise asymptotic coverage $1-\alpha$ for $\theta_0$.
\end{corollary}
These results justify using the typical variance and confidence interval
approaches used throughout EIF-based causal inference.

The two instantiations below reach Theorem~\ref{thm:main}'s conclusion by
different arguments. Instantiation~2 verifies Condition~(P) directly: it
anchors the estimator at the fixed reference partition $\tau_{0,j}$, where
the usual fixed-partition expansion for a sufficient nuisance applies
directly, and shows that the perturbation from $\tau_{0,j}$ to $\tauhat_j$
reassigns only $O_p(1)$ observations, which is enough (given only a finite
second moment on $Y$, per Remark~\ref{rem:P2-mechanism}) for the resulting
disturbance to $\sqrt n(\thetahat-\theta_0)$ to vanish. Instantiation~1
instead uses uniform asymptotic linearity over the entire fixed-grid
sufficient class: that independent argument does not select a unique anchor
and requires no analogue of (P2), because the grid-exact sufficient class
$\cS_j$ need not be a singleton. In both cases the limiting influence
function is $\psi_0$, so the variance is $V=\E[\psi_0^2]$, the \emph{same}
variance as if a sufficient partition were known in advance: learning the
thresholds, whether confined to a fixed grid or continuum-valued, costs
nothing asymptotically. Full proofs are in the companion theory document.

\subsection{Instantiation 1: fixed grid, exact sparsity}
\label{sec:instantiation1}

We are in this setting when the sufficient classes are not empty:
$\cS_\mu\ne\emptyset$, $\cS_e\ne\emptyset$. Sparsity implies that both
nuisances are representable by trees that observe the leaf budget $\Lbar$ on
the analyst's own pre-specified grid. This is the expository case: the grid
is finite by construction (\S\ref{sec:trees}), so once selection lands in
the sufficient class, topology recovery is exact. The resulting central
limit theorem is proved by uniform asymptotic linearity over the whole
sufficient class, not by verifying Condition~(P) or comparing the selected
partition with a single reference partition -- see the discussion following
Proposition~\ref{lem:uniform} below.

First, note that the doubletree nuisance functions will land in the
sufficient class with probability tending to 1:

\begin{proposition}[Selection lands in the sufficient class]
\label{lem:selection}
Suppose the observations are \iid; Assumptions~\ref{ass:causal} and
\ref{ass:construct} hold; the conditional second moments of $Y$ are uniformly
bounded; the pre-specified cutpoints and their induced positive-mass grid
atoms are fixed and finite in $n$, and optimization considers only partitions
built from those cutpoints; $\Lbar$ is fixed;
$\Lbar m_n\lesssim n$; and $\tauhat_j$ is an exact global minimizer
of~\eqref{eq:select}, with $\lambda_n\to0$. If
$\cS_e\ne\emptyset$ and $\cS_\mu\ne\emptyset$, then, for each
$j\in\{e,\mu\}$,
\[
  \Prob\bigl(\tauhat_j\in\cS_j\bigr)\longrightarrow1 .
\]
\end{proposition}

This proposition ensures that $\tauhat_j$ is consistent for the true nuisance $\tau$ for both propensities $j = e, \mu$. If we furthermore have $Y$ bounded, we can establish the rate of this selection consistency:

\begin{proposition}[Exponential selection bound]
\label{prop:selection-rate}
Under the conditions of Proposition~\ref{lem:selection}, suppose additionally
that $Y$ is bounded. Define
\[
  \Delta_j
  =
  \min_{\tau\in\cT_{\Lbar}\setminus\cS_j}
  \left\{
    R^{(j)}(\tau)
    -
    \min_{\tau'\in\cT_{\Lbar}}R^{(j)}(\tau')
  \right\}.
\]
Then $\Delta_j>0$, and for some constant $c_1>0$, depending only on the
outcome bound, $c$, and the positive leaf masses of the relevant sufficient
partitions,
\[
  \Prob\bigl(\tauhat_j\notin\cS_j\bigr)
  \le
  \card{\cT_{\Lbar}}\exp\{-c_1n\Delta_j^2\},
  \qquad j\in\{e,\mu\}.
\]
\end{proposition}

The consistency of the nuisances implies the consistency of the doubletree estimator

\medskip
\noindent\textbf{Corollary (Consistency).}
Under the conditions of Proposition~\ref{lem:selection},
\[
  \thetahat-\theta_0=o_p(1).
\]
\medskip

We now establish Theorem~\ref{thm:main}'s conclusion by an argument that
bypasses Condition~(P) rather than verifying it, unlike the
anchor-and-perturbation strategy sketched in \S\ref{sec:partition-recovery}:
because the grid-exact sufficient class $\cS_j$ need not be a singleton, we
establish uniform asymptotic linearity over the \emph{whole} class, rather
than anchoring at one fixed member of it.

\begin{proposition}[Uniform asymptotic linearity on the sufficient class]
\label{lem:uniform}
For $(\tau_e,\tau_\mu)\in\cT_{\Lbar}^2$, let
$\thetahat(\tau_e,\tau_\mu)$ denote the estimator formed from the corresponding
full-sample leaf refits. Under the conditions of Proposition~\ref{lem:selection},
\[
  \max_{(\tau_e,\tau_\mu)\in\cS_e\times\cS_\mu}
  \left|
    \sqrt n\{\thetahat(\tau_e,\tau_\mu)-\theta_0\}
    -
    \sqrt n\,\Prob_n\psi_0
  \right|
  =o_p(1).
\]
Exact global optimization and the rate of $\lambda_n$ are not required for
this uniform statement.
\end{proposition}

This proposition establishes an asymptotic equivalence between the normalized empirical mean of the influence function and the normalized error for \emph{every} estimator based on nuisances drawn from the sufficient classes. This result is a statement not yet about the particular doubletree estimator but about any estimator using sufficient nuisances.

\begin{proposition}[Fixed-grid central limit theorem by uniform linearity]
\label{prop:P-instantiation1}
Suppose the conditions of Proposition~\ref{lem:uniform} hold,
$\cS_e\ne\emptyset$, $\cS_\mu\ne\emptyset$, and $V:=\E[\psi_0^2]>0$. Then
\[
  \sqrt n(\thetahat-\theta_0)\rightsquigarrow N(0,V).
\]
This conclusion follows independently of Condition~(P): no (P1) or (P2)
verification is asserted, and in particular the argument needs no unique
reference partition, only that selection lands in the (possibly
non-singleton) sufficient class.
\end{proposition}

\begin{proof}
By Proposition~\ref{lem:selection}, $\tauhat_j\in\cS_j$ for both $j=e,\mu$
with probability tending to one. On that event, Proposition~\ref{lem:uniform}
gives
\[
  \sqrt n(\thetahat-\theta_0)
  =
  \sqrt n\,\Prob_n\psi_0+o_p(1).
\]
The classical i.i.d.\ central limit theorem and Slutsky's theorem -- the same
concluding step used for Theorem~\ref{thm:main} in the companion theory
document -- give the stated conclusion, with $V>0$ by hypothesis.
\end{proof}

\begin{corollary}[Single-outcome-tree plug-in as a special case of doubletree]
\label{cor:single-tree}
Fix the outcome tree's leaf-refit $\hat\mu=\hat\mu_{\tauhat_\mu}$ as in \S\ref{sec:trees}, fit with no propensity tree at all, and consider the treated-mean plug-in
\[
  \thetahat_{\mathrm{tie}}
  \;=\;
  \frac{1}{n_1}\sum_{i:A_i=1}\bigl\{Y_i-\hat\mu(\bX_i)\bigr\},
  \qquad
  n_1=\sum_{i=1}^n A_i.
\]
For a leaf $\ell$ of $\tauhat_\mu$ write $|\ell|_1=|\ell|-|\ell|_0$ for its treated count, and let $\hat w_{\mathrm{tie}}(\bx)=|\ell_{\tauhat_\mu}(\bx)|_1/|\ell_{\tauhat_\mu}(\bx)|_0$ be the empirical treated-to-control odds within the leaf containing $\bx$. Then $\thetahat_{\mathrm{tie}}$ solves the doubletree estimating equation $\Prob_n[\psi(\bO;\theta,\hat e,\hat\mu)]=0$ for $\hat e=\hat w_{\mathrm{tie}}/(1+\hat w_{\mathrm{tie}})$, and in fact for \emph{any} leaf-wise-constant-on-$\tauhat_\mu$ choice of $\hat e$: it \emph{is} doubletree's own two-tree estimator with the propensity tree forced equal to the outcome tree, not an alternative estimator that avoids a propensity altogether. Because the point estimate is invariant to which leaf-wise-constant $\hat e$ is used, this identity alone does not yet pin down $\thetahat_{\mathrm{tie}}$'s asymptotic variance -- that requires the specific choice $\hat e=\hat w_{\mathrm{tie}}/(1+\hat w_{\mathrm{tie}})$ and the separate argument of Lemma~\ref{lem:single-tree-linear} below.
\end{corollary}

\begin{proof}
Within any leaf $\ell$ of $\tauhat_\mu$, $\hat\mu(\bx)=\hat\mu(\ell)$ is by construction the control-arm sample mean, so $\sum_{i\in\ell}(1-A_i)\{Y_i-\hat\mu(\ell)\}=0$ identically -- for \emph{any} constant multiplying it. Hence $\Prob_n[(1-A)\hat e/(1-\hat e)\{Y-\hat\mu(\bX)\}]=\sum_\ell \hat w(\ell)\cdot n^{-1}\sum_{i\in\ell}(1-A_i)\{Y_i-\hat\mu(\ell)\}=0$ for any leaf-wise-constant $\hat e$ on $\tauhat_\mu$, with $\hat w=\hat e/(1-\hat e)$; this includes $\hat e=\hat w_{\mathrm{tie}}/(1+\hat w_{\mathrm{tie}})$ as one case among many. The doubletree estimating equation therefore reduces exactly to $\Prob_n[A\{Y-\hat\mu(\bX)-\theta\}]=0$, whose solution is $\thetahat_{\mathrm{tie}}$ by definition.
\end{proof}

Write $e_{\mathrm{tie}}(\bx)=\Prob(A=1\mid \bX\in\ell_{\tauhat_\mu}(\bx))$ for the population-level leaf-average propensity implicit in the tied construction, $w_{\mathrm{tie}}=e_{\mathrm{tie}}/(1-e_{\mathrm{tie}})$, and $\mu_{\mathrm{tie}}=\Pi^\nu_{\tauhat_\mu}\mu_0$ for the control-weighted leaf projection of $\mu_0$ defined in \S\ref{sec:trees}; write $\eta_{\mathrm{tie}}=(e_{\mathrm{tie}},\mu_{\mathrm{tie}})$, matching the $\eta=(e,\mu)$ convention of \S\ref{sec:setup}. Note $1-e_{\mathrm{tie}}(\bx)=1-\E[e_0(\bX)\mid\bX\in\ell]\ge c$ by Assumption~\ref{ass:causal}, so $\eta_{\mathrm{tie}}$ satisfies Proposition~\ref{prop:bilinear}'s hypothesis. Let $\psi_{\mathrm{tie}}=\psi(\bO;\theta_{\mathrm{tie}},\eta_{\mathrm{tie}})$ denote doubletree's score \eqref{eq:score} evaluated at these tied nuisances, where $\theta_{\mathrm{tie}}$ solves $\E[\psi_{\mathrm{tie}}]=0$.

\begin{lemma}[Asymptotic linearity of the tied estimator]
\label{lem:single-tree-linear}
Suppose $\cS_\mu\ne\emptyset$. On the event $\tauhat_\mu\in\cS_\mu$, which has probability tending to one by Proposition~\ref{lem:selection}, $\thetahat_{\mathrm{tie}}$ is asymptotically linear with influence function $\psi_{\mathrm{tie}}$:
\[
  \sqrt n(\thetahat_{\mathrm{tie}}-\theta_{\mathrm{tie}})=\sqrt n\,\Prob_n\psi_{\mathrm{tie}}+o_p(1).
\]
\end{lemma}

\begin{proof}
Since $\cS_\mu\subseteq\cT_{\Lbar}$ is a fixed, finite set, it suffices to establish the expansion for each fixed $\tau\in\cS_\mu$ and take a maximum over finitely many such expansions -- the same device used for Term~(II) in the proof of Theorem~\ref{thm:main}. At fixed $\tau$, $\thetahat_{\mathrm{tie}}=\sum_{\ell\in\tau}(|\ell|_1/n_1)\{\bar Y_1(\ell)-\hat\mu(\ell)\}$ is a smooth (ratio-of-sums) function of the finitely many leaf-wise treated counts, treated-arm means, and control-arm means, each an ordinary sample average converging at rate $n^{-1/2}$ to its population counterpart. A first-order Taylor expansion of this finite-dimensional smooth function gives asymptotic linearity at fixed $\tau$, with influence function equal to the efficient influence function of the ATT in the model that coarsens $\bX$ to its leaf $\ell_\tau(\bX)$ -- namely $\psi(\cdot;\theta_{\mathrm{tie}},e_\tau,\mu_\tau)$ for $e_\tau(\ell)=\Prob(A=1\mid\bX\in\ell)$ and $\mu_\tau(\ell)=\Pi^\nu_\tau\mu_0(\ell)$, which are exactly $e_{\mathrm{tie}}$ and $\mu_{\mathrm{tie}}$ restricted to $\tau$. Taking the maximum over the finitely many $\tau\in\cS_\mu$ preserves the $o_p(n^{-1/2})$ remainder.
\end{proof}

\begin{corollary}[Exact efficiency under saturated leaves]
\label{cor:single-tree-saturated}
If $e_0$, and not only $\mu_0$, is $\Prob$-a.s.\ constant on every leaf of $\tauhat_\mu$ -- a partition sufficient for \emph{both} nuisances simultaneously, i.e.\ $\tauhat_\mu\in\cS_e\cap\cS_\mu$ -- then $e_{\mathrm{tie}}=e_0$ and $\mu_{\mathrm{tie}}=\mu_0$ $\Prob$-a.s., so $\psi_{\mathrm{tie}}=\psi_0$ identically. Since $(\tauhat_\mu,\tauhat_\mu)\in\cS_e\times\cS_\mu$ in this case, Proposition~\ref{lem:uniform} applies directly to $\thetahat(\tauhat_\mu,\tauhat_\mu)=\thetahat_{\mathrm{tie}}$ (Corollary~\ref{cor:single-tree}), so $\sqrt n(\thetahat_{\mathrm{tie}}-\theta_0)\rightsquigarrow N(0,V)$ for $V=\E[\psi_0^2]$, the same semiparametric efficiency bound attained by doubletree's own two-tree estimator (Theorem~\ref{thm:main}), with no propensity tree required.
\end{corollary}

\begin{corollary}[A super-efficient submodel estimator under coarsening and homoskedasticity]
\label{cor:single-tree-coarsening}
Suppose only $\mu_0$ is $\Prob$-a.s.\ constant on every leaf of $\tauhat_\mu$, i.e.\ $\tauhat_\mu\in\cS_\mu$ but not necessarily $\cS_e$, so $\mu_{\mathrm{tie}}=\mu_0$ but $e_{\mathrm{tie}}$ need not equal $e_0$. Suppose in addition that Condition~(P) holds for the propensity as well (equivalently $\cS_e\ne\emptyset$ under Instantiation~1), so that the untied two-tree estimator attains $V=\E[\psi_0^2]$ by Theorem~\ref{thm:main}; the two hypotheses are compatible, and their compatibility is the point. Then $\theta_{\mathrm{tie}}=\theta_0$ \emph{exactly}, since the bilinear remainder of Proposition~\ref{prop:bilinear} vanishes when its $(\mu_0-\mu_{\mathrm{tie}})$ factor is identically zero, and by Lemma~\ref{lem:single-tree-linear}, $\thetahat_{\mathrm{tie}}$ is $\sqrt n$-consistent and asymptotically normal with variance $\Var(\psi_{\mathrm{tie}})$.

If, further, $\sigma_0^2(\bx)$ is constant on every leaf of $\tauhat_\mu$ (within-leaf homoskedasticity of the control outcome), then
\[
  \Var(\psi_{\mathrm{tie}})\ \le\ \Var(\psi_0)=V,
\]
strictly whenever $e_0$ is not $\Prob$-a.s.\ constant on some leaf with $\sigma_{0,\ell}^2>0$. The treated-arm terms of $\psi_{\mathrm{tie}}$ and $\psi_0$ coincide exactly, since $\mu_{\mathrm{tie}}=\mu_0$ and $\theta_{\mathrm{tie}}=\theta_0$; the two variances therefore differ only through the control-arm term. Doubletree's own control-arm term is $\pi^{-2}\E[e_0^2\sigma_0^2/(1-e_0)]=\pi^{-2}\E[f(e_0(\bX))\sigma_0^2(\bX)]$ for the strictly convex function $f(e)=e^2/(1-e)$ ($f''(e)=2/(1-e)^3>0$); under homoskedasticity, the tied estimator's corresponding term is $\pi^{-2}\sum_{\ell\in\tauhat_\mu}\Prob(\ell)\,\sigma_{0,\ell}^2\,f(\bar e_\ell)$ for the leaf average $\bar e_\ell=\E[e_0(\bX)\mid \bX\in\ell]$, while doubletree's own term is $\pi^{-2}\sum_\ell\Prob(\ell)\,\sigma_{0,\ell}^2\,\E[f(e_0)\mid\ell]$; Jensen's inequality applied to $f$ within each leaf gives $\E[f(e_0)\mid\ell]\ge f(\bar e_\ell)$, with strict inequality whenever $e_0$ is not $\Prob$-a.s.\ constant on $\ell$. \textbf{Heteroskedasticity breaks this ordering}: if $\sigma_0^2(\bx)$ varies within a leaf and correlates with $e_0(\bx)$ there, the per-leaf comparison is no longer a within-leaf Jensen step, and $\thetahat_{\mathrm{tie}}$ can be less efficient than the untied estimator -- indeed substantially so (a two-point leaf with $e_0\in\{0.1,0.8\}$ equally likely and $\sigma_0^2\in\{1,0\}$ respectively already gives the untied term $\approx0.006$ against the tied term $\approx0.30$, roughly a fiftyfold reversal).

This ordering does not contradict Theorem~\ref{thm:regular}: $\thetahat_{\mathrm{tie}}$ is \emph{not} a regular estimator of $\theta_0$ in the full (unrestricted) model. Under within-leaf homoskedasticity, $\psi_{\mathrm{tie}}$ is exactly the projection of $\psi_0$ onto the tangent space of the strictly smaller submodel in which $\mu_0$ is constant on the leaves of $\tauhat_\mu$, so $\Var(\psi_{\mathrm{tie}})$ is the semiparametric efficiency bound \emph{for that submodel}, not for the unrestricted model that Theorem~\ref{thm:regular} bounds; $\thetahat_{\mathrm{tie}}$ attains the submodel's (smaller) bound precisely by exploiting the restriction, which is exactly the local-alternative sensitivity that Remark~\ref{rem:single-tree-negative} shows below: moving $\mu_0$ off the submodel introduces a first-order bias, the hallmark of an estimator that is efficient within a submodel but irregular outside it. The same mechanism explains why homoskedasticity matters: it is precisely the condition under which the unweighted within-leaf mean is the efficient estimator of the leaf's control mean, so that $\mu_{\mathrm{tie}}$'s implicit weighting is optimal for the submodel; under heteroskedasticity it is not, and the projection argument no longer guarantees an improvement.
\end{corollary}

\begin{remark}[The naive standard error under-covers]
\label{rem:single-tree-se}
Because $\hat\mu$ is itself estimated from the full sample, treating it as fixed when computing $\Var(\thetahat_{\mathrm{tie}})$ omits the nonnegative control-arm variance contribution $\pi^{-2}\E[\{1-e_0(\bX)\}\hat w_{\mathrm{tie}}(\bX)^2\sigma_0^2(\bX)]$ entirely -- the naive variance is the treated-arm term alone -- and so generically under-covers. Under the hypotheses of Corollary~\ref{cor:single-tree-coarsening}, the correct asymptotic variance is $\Var(\psi_{\mathrm{tie}})$; in general it is whatever variance Lemma~\ref{lem:single-tree-linear} supplies. Either way, one never actually escapes the propensity by omitting its tree: the correct standard error recovers a propensity-like term as a byproduct of the outcome tree's own leaf structure.
\end{remark}

\begin{remark}[Distinct from a single shared-structure tree]
\label{rem:single-tree-distinct}
Corollary~\ref{cor:single-tree} and its consequences answer a different question from a single \emph{shared-structure} tree explored in earlier working notes (one tree serving both nuisances via leave-one-out algorithmic stability, following \citet{ChenSyrgkanisAustern2022}): there both nuisances are refit on the \emph{same} partition and the target is a working-model estimand, not necessarily $\theta_0$ itself. Here only the outcome is fit by a tree at all -- the propensity is never estimated as an independent object -- and the target throughout remains the true ATT $\theta_0$. These are different mechanisms answering different questions and should not be conflated.
\end{remark}

\subsection{Instantiation 2: continuum sparsity and off-grid threshold refinement}
\label{sec:instantiation2}

The fixed grid of \S\ref{sec:trees} is a convenience for the analyst, not a
property of the underlying propensity or outcome regression. This
instantiation verifies Condition~(P) when the true nuisances are exactly
representable by \emph{some} axis-aligned tree of the same fixed size
$\Lbar$, but at threshold values that need not coincide with any
pre-specified cutpoint -- so-called \emph{continuum sparsity}. The relevant
tree-fitting machinery -- a two-stage estimator that recovers the true tree
\emph{topology} on a search grid that is allowed to refine with $n$, then
refines each recovered threshold by an exact scan over the observed
covariate values -- is developed and proved in a companion paper on
continuum-threshold recovery for optimal trees; we restate here, in full,
exactly what that paper's guarantees require and deliver, since Doubletree's
own efficiency claim depends on them.

\begin{assumption}[Continuum tree sparsity, jump regularity, and local design]
\label{ass:sparsity-jump}
For each nuisance $\gamma_{0,j}$, $j\in\{e,\mu\}$, with all constants fixed
uniformly over the true internal nodes and in $n$: (S) \emph{Tree
realisability} -- $\gamma_{0,j}$ is piecewise constant on a continuum
partition realised by an axis-aligned tree with at most $\Lbar$ leaves and
fixed depth $\bar D$, with $\lceil\log_2\Lbar\rceil\le\bar D$; (J)
\emph{Genuine jumps, pooled across each split} -- every true split separates
regions whose \emph{pooled} population values differ by an amount between
fixed constants $0<\kappa_{\min}\le\kappa_{\max}<\infty$; (Local-Jump)
\emph{Jumps across every adjacent leaf pair belonging to a true split} --
every pair of leaves adjacent across a true boundary differs in population
value by between $\kappa_{\min}$ and $\kappa_{\max}$; (Reduced)
\emph{Reducedness of the true partition} -- (Local-Jump) holds for
\emph{every} geometrically adjacent pair of true leaves, whether or not the
shared face is the face of one true split -- equivalently, no two
geometrically adjacent true leaves carry the same value; this is
\emph{not} implied by minimality of the true partition (Paper A exhibits an
L-shaped counterexample) and is an independent hypothesis; (Dens)
\emph{Local connected density} -- the covariate law has a density bounded
between fixed positive constants on a fixed neighborhood of every true
splitting boundary, on supports whose relevant boxes have bounded diameter;
(Sep) \emph{Separation and interiority} -- distinct true boundaries on the
same coordinate are separated by a fixed $\varsigma>0$, and every true
boundary is at least $\varsigma$ from the support boundary; (Node)
\emph{Positive node mass} -- every true node and leaf has $\nu_j$-mass at
least a fixed $q_{\min}>0$; (Joint) \emph{Ancestor--descendant joint
density} -- covariates on distinct splitting coordinates have a jointly
bounded density near any ancestor--descendant pair of true boundaries.
\end{assumption}

Assumption~\ref{ass:sparsity-jump}'s jump and separation clauses (J), (Sep),
and (Reduced) are what rule out topological ties and give continuum
sparsity a \emph{unique} true partition, unlike the grid-exact case of
\S\ref{sec:instantiation1}.

\begin{assumption}[Two-stage solver fidelity]
\label{ass:solver-twostage}
Stage~A is an exact global minimizer of the penalized empirical risk over
the grid-restricted tree class \emph{at the true budgets} $(\Lbar,\bar D)$
-- no inflated budget or collapsed topology is used -- built on a search
grid of resolution $r_n\to\infty$ satisfying the joint rate conditions
\eqref{eq:stage-a-rates-manuscript} below. Stage~B then
performs a sequence of exact one-threshold scans, root-first, each over the
observed covariate values in a \emph{shrinking} search interval around the
corresponding Stage-A threshold, of half-width $M_n/r_n$ with $M_n\to\infty$
and $M_n/r_n\to0$; every such interval contains the corresponding true
threshold, and no other true threshold on the same coordinate, with
probability tending to one. Only Stage~B may evaluate continuous or
observed-value thresholds -- relaxing precisely this restriction, without
paying for it in the combinatorial search, is the purpose of the two-stage
procedure. Both stages use all $n$ observations, with no sample splitting;
Stage~B moves thresholds only, adding, deleting, reordering, or
re-coordinating no split; both stages use deterministic, measurable
tie-breaking; and every leaf of the refined partition contains, with
probability tending to one, at least a fixed positive fraction of $n$
observations (respectively, control observations).
\end{assumption}

For a search-grid resolution $r_n$, a penalty $\lambda_n$, and a minimum-leaf
floor $m_n$, the rate conditions this section's results require, restated
verbatim from the companion paper, are
\begin{equation}\label{eq:stage-a-rates-manuscript}
\begin{aligned}
  &\text{(R1)}&&r_n\to\infty,
    \qquad r_n=o\bigl(\sqrt{n/\log r_n}\bigr);\\
  &\text{(R2)}&&\lambda_n\to0;\\
  &\text{(R3)}&&
    \max\Bigl(\frac1{r_n},\sqrt{\frac{\log r_n}{n}}\Bigr)
    \ll\lambda_n\ll1;\\
  &\text{(R4)}&&m_n\to\infty,
    \qquad\frac{m_nr_n}{n}\to0,
    \qquad m_n\lambda_n\to\infty;\\
  &\text{(R5)}&&\frac{m_n}{\sqrt n}\to\infty.
\end{aligned}
\end{equation}
Condition (R3) requires $\lambda_n\gg1/r_n$, since the penalty must
dominate the Stage-A grid-approximation error -- the reverse of the
fixed-grid regime's $\lambda_n\Lbar<\Delta_j/2$ of
\S\ref{sec:partition-recovery}; the two statements concern different
regimes and no equivalence between their proof mechanisms is asserted.

\begin{proposition}[Topology recovery, restated from the companion paper]
\label{prop:boundary-recovery}
Under Assumptions~\ref{ass:causal}, \ref{ass:construct},
\ref{ass:sparsity-jump}, and~\ref{ass:solver-twostage}'s Stage-A clause, and
a search-grid resolution $r_n$ satisfying the joint rate conditions
\eqref{eq:stage-a-rates-manuscript}, every Stage-A penalised global
minimiser satisfies, with probability tending to one: its number of splits
equals the true partition's; it structurally matches the true partition,
node for node, in the sense that every Stage-A node corresponds to a true
node at the same relative position with the same splitting coordinate; and,
under that correspondence, every Stage-A threshold is within $O(r_n^{-1})$
of its matched true threshold. The result uses the actual budgets
$(\Lbar,\bar D)$; no inflated budget or collapse map appears.
\end{proposition}

\begin{proposition}[Threshold superconsistency, restated from the companion paper]
\label{prop:twostage-superconsistency}
Under the full hypotheses of Proposition~\ref{prop:boundary-recovery} for
both nuisances, Assumption~\ref{ass:sparsity-jump}(Joint), all clauses of
Assumption~\ref{ass:solver-twostage}, and the Stage-B interval scale
conditions $M_n\to\infty$, $M_n/r_n\to0$, every Stage-B-refined threshold
satisfies $\hat c_{j,k}-c_{0,j,k}=O_p(n^{-1})$, simultaneously over
$k=1,\ldots,K_j\le\Lbar-1$, for each $j\in\{e,\mu\}$.
\end{proposition}

\begin{proposition}[Condition~(P) holds under continuum sparsity]
\label{prop:P-instantiation2}
Under the full hypotheses of Propositions~\ref{prop:boundary-recovery} and
\ref{prop:twostage-superconsistency} for both nuisances, Condition~(P) holds
for both $j=e,\mu$: (P0) is Assumption~\ref{ass:sparsity-jump}(S), with
positive leaf mass from (Node); (P1) is
Proposition~\ref{prop:boundary-recovery}'s structural-matching conclusion,
read at the level of the induced leaf correspondence (topology and matched
split coordinates; threshold values may differ, and equality of the induced
leaf sets is not asserted); and (P2) follows from
Proposition~\ref{prop:twostage-superconsistency}'s $O_p(n^{-1})$ rate: by
the same order-statistic argument underlying Remark~\ref{rem:P2-mechanism},
this threshold rate, together with the local density bounds of
Assumption~\ref{ass:sparsity-jump}(Dens), implies $\card{S_{j,n}}=O_p(1)$.
Theorem~\ref{thm:main} therefore applies with the \emph{same} variance
$V=\E[\psi_0^2]$ as Instantiation~1, via the anchor-and-perturbation
argument of \S\ref{sec:partition-recovery}.
\end{proposition}

\begin{remark}[The grid-exactness cost is retired]
\label{rem:apology-retired}
Fixing the grid in advance no longer forces the truth to align with the
analyst's particular choice of cutpoints: Proposition~\ref{prop:P-instantiation2}
shows that exact recovery -- topology and, up to an $O_p(n^{-1})$ threshold
error, location -- is available for any true boundary in the continuum,
not merely for boundaries the analyst happened to place a cutpoint on. What
remains required is \emph{structural} sparsity in the wider sense of
Assumption~\ref{ass:sparsity-jump}: the truth must still be exactly
piecewise constant on some fixed-size tree, with genuine, separated jumps.
Refining the analyst's search grid relaxes the first cost (grid alignment)
but not the second (exact tree-shaped structure); \S\ref{sec:instantiation3}
addresses that second cost directly.
\end{remark}

\begin{remark}[Outside exact leaf-measurability, the second factor is no longer independently controllable]
\label{rem:single-tree-negative}
Corollary~\ref{cor:single-tree-coarsening}'s exactness $\theta_{\mathrm{tie}}=\theta_0$ required $\mu_0$ to be exactly $\Prob$-a.s.\ constant on every leaf of $\tauhat_\mu$, i.e.\ $\tauhat_\mu\in\cS_\mu$. In this instantiation, or under genuinely misspecified sparsity, that generically fails, and tying the propensity to the outcome tree's own partition does not lose the bilinear (product) structure of Proposition~\ref{prop:bilinear} -- but it does force the second factor to track the \emph{same} partition as the first, removing the analyst's ability to control it independently. Applying Proposition~\ref{prop:bilinear} at $\eta_{\mathrm{tie}}=(e_{\mathrm{tie}},\mu_{\mathrm{tie}})$ and using $\{1-e_0(\bX)\}w_0(\bX)=e_0(\bX)$,
\[
  \theta(\eta_{\mathrm{tie}})-\theta_0
  \;=\;
  \frac1\pi\,\E\bigl[\{1-e_0(\bX)\}\{w_0(\bX)-w_{\mathrm{tie}}(\bX)\}\{\mu_0(\bX)-\mu_{\mathrm{tie}}(\bX)\}\bigr]
  \;=\;
  \frac1\pi\,\E\bigl[e_0(\bX)\{\mu_0(\bX)-\mu_{\mathrm{tie}}(\bX)\}\bigr],
\]
since $\E[\{1-e_0(\bX)\}w_{\mathrm{tie}}(\bX)\{\mu_0(\bX)-\mu_{\mathrm{tie}}(\bX)\}\mid \bX\in\ell]=w_{\mathrm{tie}}(\ell)\,\E[\{1-e_0(\bX)\}\{\mu_0(\bX)-\mu_{\mathrm{tie}}(\ell)\}\mid \bX\in\ell]=0$ leaf by leaf, $\mu_{\mathrm{tie}}(\ell)$ being exactly the $\{1-e_0\}$-weighted leaf mean of $\mu_0$. Writing $m_\ell=\E[\mu_0(\bX)\mid\bX\in\ell]=\Pi_\tau\mu_0(\ell)$ for the \emph{unweighted} leaf mean and $\bar e_\ell=e_{\mathrm{tie}}(\ell)$, the identity $\E[e_0(\mu_0-\mu_{\mathrm{tie}})\mid\ell]=m_\ell-\mu_{\mathrm{tie}}(\ell)=\Cov(e_0,\mu_0\mid\ell)/(1-\bar e_\ell)$ recasts the bias as a sum of within-leaf covariances,
\[
  \theta(\eta_{\mathrm{tie}})-\theta_0
  \;=\;
  \frac1\pi\sum_{\ell\in\tauhat_\mu}\Prob(\ell)\,\frac{\Cov(e_0,\mu_0\mid\ell)}{1-\bar e_\ell},
\]
and Cauchy--Schwarz within each leaf, followed by Cauchy--Schwarz across leaves and $1-\bar e_\ell\ge c$ (Assumption~\ref{ass:causal}), gives the explicit bound
\[
  \bigl|\theta(\eta_{\mathrm{tie}})-\theta_0\bigr|
  \;\le\;
  \frac{1}{\pi c}\,\|e_0-\Pi_\tau e_0\|_2\,\|\mu_0-\Pi_\tau\mu_0\|_2 .
\]
The product-of-two-errors structure is therefore \emph{not} lost -- it is Lemma~\ref{lem:biasbound}'s bound, restated with both factors evaluated at the \emph{same} partition $\tauhat_\mu$ rather than at two independently chosen ones. That loss of independent control, not a loss of bilinearity, is the precise reason two \emph{independent} trees -- not one tree, not one shared-structure tree -- are needed outside grid-exact sparsity: with two trees, one nuisance's leaf-measurability alone already annihilates the bias (Proposition~\ref{prop:bilinear}); with one tree serving both roles, the bias survives whenever \emph{either} nuisance has residual within-leaf heterogeneity on that shared partition, and vanishes only if the product $\|e_0-\Pi_\tau e_0\|_2\|\mu_0-\Pi_\tau\mu_0\|_2$ itself is $o(n^{-1/2})$ -- a condition attainable with both factors merely $O_p(n^{-1/4})$, not requiring the change-point-style single-factor rate that a narrower reading of the bound might suggest.
\end{remark}

\subsection{When partition recovery fails: the anchor interval}
\label{sec:honest-manuscript}
Condition~(P) holding delivers efficiency (\S\ref{sec:partition-recovery}--\S\ref{sec:instantiation2}); this section addresses what happens when it may not hold. We are in the non-sparse setting when at least one sufficient class is empty: $\cS_e = \emptyset$ or $\cS_\mu = \emptyset$. Here, at least one nuisance is not exactly representable by a tree that observes the leaf budget $\Lbar$, so we should expect some bias in $\thetahat$ relative to $\theta_0$. The interval constructed below is, in spirit, an honest, approximation-aware confidence interval in the tradition of \citet{armstrong2018}: rather than assuming the bias away or estimating it directly, it is bounded and absorbed into the interval's width.

First, note that the doubletree estimator admits an exact decomposition of this bias:

\begin{proposition}[Exact bilinear remainder]
\label{prop:bilinear}
Under Assumption~\ref{ass:causal}, for any
$\eta=(e,\mu)$ satisfying $1-e(\bx)\ge c$ for every $\bx\in\cX$,
\begin{equation}
\label{eq:bilinear}
  \E\bigl[\psi(\bO;\theta_0,\eta)\bigr]
  =
  \frac{1}{\pi}
  \E\left[
    \{1-e_0(\bX)\}
    \{w_0(\bX)-w(\bX)\}
    \{\mu_0(\bX)-\mu(\bX)\}
  \right].
\end{equation}
\end{proposition}

This decomposition shows that the bias of $\thetahat$ is exactly a bilinear remainder in the errors of the two fitted nuisances -- the same mechanism that gives doubly robust estimators their name. Bounding each nuisance's error in the appropriate norm gives an explicit bound on the size of that bias:

\begin{lemma}[Bias bound in the control-weighted norm]
\label{lem:biasbound}
Under Assumptions~\ref{ass:causal} and~\ref{ass:construct}, define the
realized nuisance errors
\[
  D_w=\|\hat w-w_0\|_{2,w},
  \qquad
  D_\mu=\|\hat\mu-\mu_0\|_{2,w}.
\]
Then the population estimating-equation bias evaluated at the fitted
nuisances satisfies
\[
  \left|
    \left.
    \E\bigl[\psi(\bO;\theta_0,\eta)\bigr]
    \right|_{\eta=\hat\eta}
  \right|
  \le
  \frac{1}{\pi}D_wD_\mu .
\]
\end{lemma}

This bound shows that $|\thetahat - \theta_0|$ is controlled, up to constants, by the product of the two realized nuisance errors $D_w D_\mu$ -- not directly by the population approximation errors $\delta_e, \delta_\mu$, which measure only the best achievable bias at budget $\Lbar$ and need not equal what the fitted trees actually attain. When $D_w D_\mu$ is $o(n^{-1/2})$, $\thetahat$ remains asymptotically efficient even outside the sparse setting; but in general we cannot guarantee this rate, so we cannot guarantee that the usual normal-theory confidence interval is valid.

We therefore need an inferential strategy that is valid regardless of $D_w D_\mu$. The key idea is to rely on a valid, non-interpretable \emph{anchor} estimate, $\thetatilde$: a second estimator that is $\sqrt n$-consistent and asymptotically normal without requiring sparsity. The anchor might be what you would choose outside of high-stakes settings requiring intepretability. In this case, it will help us obtain inference for our interpretable estimate. The anchor need not use optimal trees or cross-fitting, it need only be $\sqrt n$-consistent, but we show in the appendix that a cross-fit estimator using optimal trees can play this role under suitable regularity conditions.

\begin{assumption}[Anchor nuisance rate]\label{ass:rate}
When the anchor is constructed from $K$ fold-wise nuisance fits, with $K$
fixed, let $\widetilde e^{(-k)}$ and $\widetilde\mu^{(-k)}$ denote the
propensity and outcome fits trained without fold $k$. These fits, which need
not use optimal trees, satisfy
\[
  \max_{1\le k\le K}
  \twonorm{\widetilde e^{(-k)}-e_0}\,
  \bigl\|\widetilde\mu^{(-k)}-\mu_0\bigr\|_{2,w}
  =
  o_p(n^{-1/2}).
\]
An anchor constructed by another method may instead be used if it directly
satisfies the $\sqrt n$-consistency and asymptotic normality conditions stated
in Theorem~\ref{thm:anchor}.
\end{assumption}

Assumption~\ref{ass:rate} concerns the anchor rather than the displayed
doubletree nuisances. It is not made automatic merely by choosing the
saturated grid budget $\Lbar=M$: when $\cX$ is general, saturation separates
all grid atoms but cannot force the nuisance functions to be constant within
those atoms. With a valid anchor $\thetatilde$ in hand, we can construct an
interval centered at the doubletree estimate itself, widened just enough to
remain valid regardless of the size of the bias:

\begin{theorem}[Anchor interval]
\label{thm:anchor}
Suppose
\[
  \sqrt n(\thetatilde-\theta_0)
  \rightsquigarrow N(0,\widetilde V),
  \qquad
  \widetilde V>0,
  \qquad
  \widetilde\sigma^2\to_p\widetilde V,
\]
and define $\widehat\delta=\thetahat-\thetatilde$. Then
\[
  C_n
  =
  \left[
    \thetahat
    \pm
    \left\{
      \abs{\widehat\delta}
      +
      z_{1-\alpha/2}\frac{\widetilde\sigma}{\sqrt n}
    \right\}
  \right]
\]
satisfies
\[
  \liminf_{n\to\infty}\Prob(\theta_0\in C_n)\ge1-\alpha
\]
on $\{\pihat>0\}$, whose probability tends to one. No rate or smoothness
condition on the doubletree nuisances is required.
\end{theorem}

This interval retains $\thetahat$, and hence the two displayed trees, as the point estimate; only the width, not the point estimate, absorbs the unknown bias. The following corollary characterizes that width explicitly:

\begin{corollary}[Width]
\label{cor:width}
Suppose the conditions of Theorem~\ref{thm:anchor} hold, the doubletree
nuisances satisfy the moment, construction, and fixed-budget conditions above,
and the anchor is $\sqrt n$-efficient. With $D_w, D_\mu$ as in
Lemma~\ref{lem:biasbound}, we have
\[
  \abs{\widehat\delta}
  =
  O_p(D_wD_\mu+n^{-1/2}),
  \qquad
  \operatorname{width}(C_n)
  =
  O_p\!\left(\max\{n^{-1/2},D_wD_\mu\}\right).
\]
In particular, the interval has width $O_p(n^{-1/2})$ whenever
$D_wD_\mu=O_p(n^{-1/2})$.
\end{corollary}

This corollary shows that the interval contracts at the usual parametric rate whenever $D_w D_\mu = O(n^{-1/2})$, and otherwise plateaus at a width set by the product of the realized nuisance errors -- an explicit price for insisting on a small, displayable tree when the truth does not fall in the interpretable class.

Since $\delta_e, \delta_\mu$ are unknown in practice, it is useful to have a diagnostic for whether the sparse setting plausibly holds at all:

\begin{proposition}[Fidelity diagnostic]
\label{prop:spectest}
Under the conditions of Theorem~\ref{thm:main}, together with a rate
condition on the anchor's fold-wise nuisance errors,
\[
  \sqrt n\,\widehat\delta=o_p(1).
\]
Consequently, for any
$\widehat{\mathrm{SE}}\asymp n^{-1/2}$, the test that rejects sparsity when
\[
  \abs{\widehat\delta}
  >
  z_{1-\alpha/2}\widehat{\mathrm{SE}}
\]
has size tending to zero. If, under an alternative, each population nuisance
risk has a unique minimizer separated from its competitors by a positive
margin and $\hat\eta\to_p\eta_*$, the test is consistent whenever
\[
  b:=\E\bigl[\psi(\bO;\theta_0,\eta_*)\bigr]\ne0 .
\]
\end{proposition}

Together, these results mean that a user need never choose, in advance, between reporting an interpretable estimator and reporting a valid one: the same two trees are displayed and used as the point estimate whether or not sparsity holds, and the diagnostic together with the anchor interval's width tell the user how much confidence to place in efficiency versus mere validity.

\subsection{Efficiency without tree-shaped truth: refining the grid}
\label{sec:instantiation3}

The anchor interval of \S\ref{sec:honest-manuscript} requires nothing of the
true nuisances, but it only promises \emph{validity}: its width contracts at
the parametric rate only if the realized bias $D_wD_\mu$ happens to be
$o(n^{-1/2})$, which cannot be guaranteed when the truth is not exactly
tree-shaped. This section gives a different answer, available to an analyst
who is willing to make a design choice rather than a structural assumption:
let the discretization grid used to \emph{construct} the interpretable tree
refine as $n$ grows, while keeping the leaf budget $\Lbar$ fixed, so the
displayed tree remains exactly as small and auditable as before.

Formally, suppose the analyst's covariate $X_n=Q_n(Z)$ is a discretization,
with mesh width $h_n\to0$, of an underlying continuum covariate $Z$, and
the true propensity and outcome regressions are general functions of $Z$ --
not assumed piecewise constant on any tree, so this is \emph{not} an
instance of Condition~(P): no true topology is being recovered, because the
truth need not have one. A companion document develops the inferential
theory for this regime in full: a population-level approximation lemma
showing the grid-alignment error between the continuum truth and the best
tree the analyst's grid can represent vanishes at rate $O(\sqrt{h_n})$, so
exact leaf-constancy is required only in the limit, never of any particular
finite $n$'s grid; and, building on a genuine oracle inequality that
replaces exact selection consistency (which cannot survive grid refinement
under any uniformity assumption), three results: a cross-fitted estimator
that is $\sqrt n$-consistent, asymptotically normal, and \emph{efficient}
whenever $h_n=o(n^{-1/2})$; an anchor interval whose width contracts at a
stated, possibly slower, rate; and, under an additional joint
empirical-process condition, a full-sample estimator with the same
$\sqrt n$-consistency, asymptotic normality, and local efficiency.

This is not a strictly stronger result than Instantiation~2's
(\S\ref{sec:instantiation2}): it abandons the institutionally motivated,
pre-specified cutpoints (age brackets, income bands, clinical thresholds)
that originally motivated fixing the grid, in exchange for not needing the
truth to be exactly piecewise constant on any tree. It therefore answers a
different analyst's objection -- one who wants a small, displayable tree but
is skeptical that any real propensity or outcome regression is genuinely a
step function -- while \S\ref{sec:honest-manuscript}'s anchor interval
remains the option that requires no design choice and no smoothness at all.

\subsection{Discussion: is a fixed candidate class basically a parametric model?}
\label{sec:finite-discussion}

A natural objection to fixing the candidate class $\cT_{\Lbar}$ in advance is that the true splits are not known ahead of time, so restricting the search to a pre-specified grid looks like trading nonparametric flexibility for an essentially parametric hypothesis space. The precise cost is narrower than that, and it is worth separating two different things a fixed, finite candidate class restricts: the analyst's \emph{search}, and the statistical \emph{model} in which $\theta_0$ is defined and efficiency is assessed. Only the former is restricted here.

\begin{theorem}[Regularity and efficiency]
\label{thm:regular}
Under the conditions of Theorem~\ref{thm:main} and a mild interiority
condition on $e_0$ (bounded away from both $0$ and $1-c$ on the interior of
its support), $\thetahat$ is regular and locally asymptotically efficient
for $\theta_0$ in the full nonparametric model: no estimator that is regular
along every parametric submodel through the true law can improve on the
asymptotic variance $V$.
\end{theorem}

\begin{proposition}[The submodel bound is not exploited]
\label{prop:sparse-bound}
Fix $\tau\in\cS_e$ and let $\cP_\tau$ denote the much smaller submodel in
which $e_0$ is assumed to be exactly constant on the leaves of $\tau$. The
semiparametric efficiency bound for $\theta_0$ in $\cP_\tau$ is $V_\tau\le
V$, strictly smaller than $V$ whenever the treated--control contrast varies
within a leaf of $\tau$: exact knowledge of the tree structure would
genuinely lower the achievable variance. By Theorem~\ref{thm:regular}, the
doubletree estimator does not exploit this -- it remains regular and attains
$V$, the bound for the full model, not $V_\tau$.
\end{proposition}

A genuinely parametric restriction on $e_0$ and $\mu_0$ would be cashed in as
a lower variance; structural sparsity is not. That is the sense in which
fixing the candidate class is not the objection it first appears to be: it
buys interpretability, and the estimator declines the resulting information
rather than exploiting it.

What \emph{does} restrict the data-generating process is structural sparsity
itself (\S\ref{sec:instantiation1}) -- that the true nuisances be
\emph{exactly} representable by some tree within the budget $\Lbar$. We do
not minimize this cost in general: Theorem~\ref{thm:main} requires the truth
to align exactly, not merely approximately, with some tree in the relevant
class. Two of the routes above already relax part of this cost, however, and
it is worth being precise about which part each relaxes. Fixing the grid in
advance -- as opposed to the truth being exactly tree-shaped -- is not itself
the binding restriction: \S\ref{sec:instantiation2} shows the same efficiency
holds when the true boundaries lie anywhere in the continuum, not merely at
the analyst's pre-specified cutpoints, citing the companion paper's
topology-recovery and off-grid threshold-refinement guarantees. What remains
required there is that the truth be \emph{exactly} piecewise constant on
some fixed-size tree with genuine, separated jumps. \S\ref{sec:instantiation3}
addresses that second, harder cost directly, at the price of no longer fixing
the analyst's institutionally motivated cutpoints at all.

The reason neither relaxation is load-bearing for the paper's inferential
guarantees is Theorem~\ref{thm:anchor}: validity does not require structural
sparsity to hold, and it does not even require the anchor to be built from the
same interpretable machinery. The anchor need only be some valid, $\sqrt
n$-consistent estimator (\S\ref{sec:honest-manuscript}); it could be a standard AIPW estimator with
flexible black-box nuisances entirely unrelated to the displayed trees, and
the anchor interval's coverage guarantee is unaffected. The question ``what
if the fixed candidate class is wrong?'' therefore has an answer that does
not depend on the fixed candidate class at all: report the same two trees,
and let the anchor -- built however one likes -- absorb the risk. What is
lost when sparsity fails is efficiency and a point estimate guaranteed to be
unbiased, not validity, and the loss is explicit and quantified, not hidden,
in Corollary~\ref{cor:width}'s width.

\subsection{Relation to penalized linear methods}
\label{sec:lasso-comparison}

Lasso-based nuisance estimation provides a useful comparison. In the orthogonal and debiased-inference literature, nuisance functions are also estimated by penalized empirical risk minimization, typically by minimizing squared loss plus an $\ell_1$ penalty \citep{belloni2014,farrell2015,chernozhukov2018}. Importantly, validity is not made conditional on exact support recovery or model-selection consistency. This separation avoids the nonuniformity associated with inference after selecting a model \citep{leeb2005}. Doubletree uses the same two-tier resolution for a tree class: it provides an efficient point estimate when structural sparsity holds, while the anchor interval remains valid when it may not.

The assumptions that purchase fast nuisance rates differ. Lasso analyses typically impose a restricted-eigenvalue or compatibility condition on the design. For doubletree, the fast selection rate in Proposition~\ref{prop:selection-rate} comes from the positive population-risk margin created by exact structural sparsity: every candidate outside the sufficient class has strictly positive excess risk, and the finite candidate class then yields a fixed gap. Finiteness alone is not enough; without the margin, a union bound over the finite class gives an excess-risk rate that is too slow for efficient inference. Doubletree's route to a fast rate therefore requires exact, rather than approximate, structural sparsity and is not a free alternative to the compatibility conditions used for Lasso.

The corresponding representational trade is also two-sided. Trees represent hierarchical or region-varying structure economically because the covariates that matter may differ across regions of the covariate space, but genuinely additive structure can require a leaf for every configuration of the active covariates. A fixed-basis linear or Lasso model has the mirror-image trade: additive structure is economical, whereas region-varying effects and interactions must be included explicitly. The relevant assumptions therefore depend on the structure of the true nuisance functions; neither representation is uniformly preferable.

\appendix

% Bibliography
\bibliographystyle{apalike}
\bibliography{refs,references}

\end{document}
